\section{成长与教育：从奥数到 OI}

\subsection{数学天赋与自驱学习}

翁家翌从小学一年级开始学奥数，并很快展现出数学方面的天赋。他描述自己的特点：\textbf{学新东西慢，但一旦理解后使用极快}。他用知识树的比喻解释这一现象——别人需要从根部沿着树枝逐层推导到结论，而他会直接建立一个 shortcut（捷径），跳过中间推导直接到达答案。

\begin{importantbox}{学习方法论：慢学习 + 快应用}
翁家翌自述需要花别人两到三倍的时间来学习新东西，读代码也比别人慢。但他的应对策略是：\textbf{提前学}。初二就学完高中数学，初三开始学微积分。他将这种行为称为"投资未来"——这一理念贯穿了他整个职业生涯。
\end{importantbox}

这种"投资未来"的思维并非来自父母的要求，而是完全自发的。他的核心逻辑是：与其把时间浪费在当前的刷题上，不如学一些对未来有用的东西，后续收益会更多。值得注意的是，这种长期主义思维在他初中时期就已经形成。

翁家翌还分享了一个关于学习方法的细节：小时候背课文，他会在睡觉之前想尽所有方法，磕磕巴巴地把全文背出来，哪怕有很多停顿也要坚持背完。然后睡一觉，第二天醒来发现倒背如流。他把这个过程比作 System 1 和 System 2 的关系——用时髦话说，口算题是 System 1 直接过，而理解新知识需要 System 2 慢慢构建，但一旦构建完成就变成了 System 1 的反射。

\subsection{编程启蒙与 OI 竞赛}

翁家翌从初一开始接触编程，起因是学校的编程兴趣班。高中时因为升学压力，他开始参加信息学竞赛（OI）。他坦言作为非北京生源，不搞竞赛想进清华"难如登天"。

在 OI 的道路上，他经历了波折：
\begin{itemize}
    \item 高一时省选几乎不会做题
    \item 凭借一道最小双元覆盖题拿下全场最高分，勉强进入福建省队
    \item 国赛（NOI）时是福建省队\textbf{倒数第一}，拿了铜牌
    \item 放弃了更保险的上海交大本一线录取，选择了清华降60分的条件录取
\end{itemize}

\begin{knowledgebox}{OI 竞赛升学路径（当时）}
竞赛升学的典型路径：NOIP（提高组省赛） $\rightarrow$ 省选 $\rightarrow$ 清华/北大夏令营（可获降分或保送条件） $\rightarrow$ NOI国赛。金牌可直接保送，银牌以上可获本一线录取条件，降60分则需要高考成绩配合。
\end{knowledgebox}

翁家翌在高三期间偷偷继续钻研 OI，甚至练成了在 iPad Safari 上直接裸打代码并提交的技能——没有编译器，没有代码编辑器。他认为这段经历极大锻炼了对整个程序逻辑的完整认知能力和快速定位 bug 的反应能力。

他还热衷于\textbf{常数优化}——在算法时间复杂度相同的情况下，优化代码运行的常数系数和代码长度。OI 评测系统会按照测试点的运行时间排序，跑得最快的排第一；如果运行时间相同，则按代码长度排序。翁家翌会同时 optimize 这两个指标——"虽然没什么用，但是很有意思"。这种对代码极致优化的追求，为后来在 OpenAI 搭建 RL Infra 埋下了伏笔。

\subsection{高考抉择：在不确定性中押注}

翁家翌在 NOI 拿了铜牌后，面临一个艰难的抉择：清华夏令营给出的条件是\textbf{降60分}（高考成绩加60分，过线即录取），但附带一个条件——如果 NOI 能进前150名（银牌线），则直接本一线录取。他最终没有拿到银牌。

另一个选项是更保险的上海交大本一线录取。翁家翌回忆说，当时高二下半年没搞文化课，不知道高考能考多少分，有学长把60分加分全部用完了，心里很害怕。但在家人的鼓励下，他选择了风险更高但天花板更高的清华降60分。

\begin{warningbox}{面对不确定性的决策模式}
这是翁家翌第一次面对重大的不确定性抉择。他的策略是：在害怕和不确定中，选择那个\textbf{上限更高}的选项。这个决策模式后来在他的职业生涯中反复出现——选 OpenAI 而非更稳妥的大厂，做开源而非发论文，都遵循同样的逻辑。
\end{warningbox}

\subsection{本章小结}

翁家翌的成长经历揭示了两个核心特质：一是\textbf{正反馈驱动的自我强化}——在擅长的领域不断获得正反馈，形成内生兴趣；二是\textbf{长期投资思维}——宁可牺牲短期效率，也要为未来积累优势。这两个特质后来深刻影响了他在清华、CMU 和 OpenAI 的每一个选择。


